\documentclass{article}

\usepackage[utf8]{inputenc}
\usepackage[T1]{fontenc}
\usepackage{helvet}
\usepackage{geometry}
\usepackage{xcolor}
\usepackage{eso-pic}
\usepackage{fancyhdr}
\usepackage{parskip}
\usepackage{microtype}
\usepackage[english, ukrainian]{babel}
\usepackage{url}
\usepackage{amsmath}
\usepackage{amssymb}
\usepackage{amsthm}
\usepackage{longtable}
\usepackage{booktabs}
\usepackage{hyperref}
\usepackage{titlesec}

\definecolor{nato}{HTML}{293757}
\definecolor{footergray}{gray}{0.65}

\geometry{a4paper,left=3cm,right=3cm,top=3cm,bottom=3cm}
\renewcommand{\familydefault}{\sfdefault}
\setcounter{secnumdepth}{3}

\titleformat{\section}
  {\color{nato}\Huge\bfseries}{\thesection.}{1em}{}
\titlespacing*{\section}{0pt}{30pt}{12pt}

\titleformat{\subsection}
  {\color{nato}\LARGE\bfseries}{\thesubsection.}{1em}{}
\titlespacing*{\subsection}{0pt}{20pt}{8pt}

\titleformat{\subsubsection}
  {\color{nato}\Large\bfseries}{\thesubsubsection.}{1em}{}
\titlespacing*{\subsubsection}{0pt}{10pt}{6pt}

\fancyhf{}
\fancyhead[R]{\small\color{footergray} 31 березня 2026}
\fancyfoot[L]{\small\color{footergray} Технічні вимоги ЄСІКС. III. Інфраструктура документообігу}

\fancyfoot[R]{\small\color{footergray}\thepage}

\newcommand\HUGE[1]{\fontsize{40}{40}\selectfont #1}

\renewcommand{\headrulewidth}{0pt}
\renewcommand{\footrulewidth}{0.4pt}
\renewcommand{\footrule}{\color{footergray}\hrule width \textwidth height 0.4pt}

\begin{document}

\pagestyle{fancy}

\begin{titlepage}
  \AddToShipoutPictureBG*{\AtPageLowerLeft{\color{nato}\rule{\paperwidth}{\paperheight}}}
  \color{white}\thispagestyle{empty}\vspace*{4cm}\centering
  {\Huge  \textbf{Технічні вимоги ЄСІКС} \par} \vspace{2cm}
  {\HUGE  \textbf{III. Документообіг} \par} \vspace{2.5cm}
  {\LARGE \textbf{III-1. Шаблони документів (DOCX, XLSX)} \par}
  {\LARGE \textbf{III-2. Генерація документів (PDF)} \par}
  {\LARGE \textbf{III-3. Генерація бар- та штрих-кодів (QR)} \par}
  {\LARGE \textbf{III-4. Розпізнавання Tesseract (OCR)} \par}
  {\LARGE \textbf{III-5. Текстуально-голосові трансформації (S2T, T2S)} \par}
  {\LARGE \textbf{III-6. Система сканування (SCAN)} \par}
  {\LARGE \textbf{III-7. Редактор справ консиліуму (CRDT)} \par} \vspace{2.5cm}
  \vfill
  {\large 2026 \copyright\ Інформаційні Судові Системи \par}
\end{titlepage}

\newpage

\begin{center}
  {\Huge\color{nato}\bfseries Зміст\par}
\end{center}
\vspace{40pt}
\tableofcontents

\begin{abstract}
ТЕХНІЧНІ ВИМОГИ
НА СТВОРЕННЯ ЗАСОБУ ІНФОРМАТИЗАЦІЇ - КОМПОНЕНТІВ ІНФРАСТРУКТУРИ ДОКУМЕНТООБІГУ
ЄДИНОЇ СУДОВОЇ ІНФОРМАЦІЙНО-КОМУНІКАЦІЙНОЇ СИСТЕМИ.
\end{abstract}

\newpage
\section*{Вступна частина}

\subsection*{Умовні скорочення та визначення}

Терміни та скорочення, що використовуються в документі:

\begin{longtable}{p{4cm}p{10cm}}
\toprule
\textbf{Термін} & \textbf{Значення} \\
\midrule
API & Автоматизовані програмні інтеграції, відомі також як “прикладний інтерфейс програмування – application programming interface” \\
БД & База Даних — Організована колекція даних, яка зберігається та управляється з використанням комп'ютерної системи; забезпечує зберігання, організацію та доступ до даних з метою ефективного використання та обробки \\
БП & Бізнес Процес — Сукупність дій необхідних для здійснення людиною або системою, в результаті яких відбувається досягнення бажаного результату \\
КЕП & Кваліфікований електронний підпис (X.509) \\
РМК & Реєстраційно-моніторингова картка \\
АРМ & Автоматизоване робоче місце користувача \\
ЄСІКС & Єдина судова інформаційно-комунікаційна система \\
BPMN & Business Process Model and Notation (ISO 19510) — стандарт моделювання та оркестрації бізнес-процесів \\
DOCX & Формат шаблонів документів Microsoft Word (Open XML) — стандартизований формат для створення та редагування текстових документів у системі документообігу \\
XLSX & Формат шаблонів електронних таблиць Microsoft Excel (Open XML) — стандартизований формат для створення та редагування табличних даних у системі документообігу \\
PDF & Portable Document Format — універсальний формат для представлення, зберігання та архівації документів незалежно від програмного забезпечення та платформи \\
QR & Quick Response code — двовимірний штрих-код для швидкого сканування, ідентифікації документів та доступу до інформації \\
OCR & Optical Character Recognition (Tesseract) — технологія автоматичного розпізнавання та перетворення тексту зі сканованих зображень і документів у редагований текст \\
S2T & Speech-to-Text — перетворення мовлення (голосу) в текст для текстово-голосових трансформацій \\
T2S & Text-to-Speech — перетворення тексту в синтезоване мовлення (голос) для текстово-голосових трансформацій \\
CRDT & Conflict-free Replicated Data Type (EDIT) — багатокористувацький суддівський редактор справ консиліуму з підтримкою одночасної роботи без конфліктів редагування \\
\bottomrule
\end{longtable}

\newpage
\subsection*{Загальні положення}

Технічні вимоги на основі «Концепції ЄСІКС» від 2024 року \cite{esiks} і авторських вимог до твору
на ``ERP/1: Документи'' \footnote{\url{https://crm.erp.uno}}.

\subsection*{Призначення та цілі впровадження}
З метою забезпечення узгодженої, безпечної та ефективної роботи ЄСІКС у сфері документообігу,
а також унеможливлення фрагментації підходів до створення, обробки, генерації та управління документами,
в межах ЄСІКС є необхідним створення та впровадження централізованих спільних компонентів, зокрема:
\begin{enumerate}
\item \textbf{Шаблони документів (DOCX, XLSX)}: централізоване створення, зберігання, версіонування та керування шаблонами документів.
\item \textbf{Генерація документів (PDF)}: автоматизована генерація, конвертація та експорт документів у PDF-форматі.
\item \textbf{Генерація та накладання бар- та штрих-кодів (QR)}: створення, інтеграція та автоматичне накладання баркодів, штрих-кодів і QR-кодів на документи.
\item \textbf{Розпізнавання Tesseract (OCR)}: автоматичне розпізнавання та перетворення тексту зі сканованих зображень і документів.
\item \textbf{Текстуально-голосові трансформації (S2T, T2S)}: підтримка двосторонніх перетворень тексту в мовлення та мовлення в текст.
\item \textbf{Багатокористувацький суддівський редактор справ консиліуму (CRDT)}: централізований редактор для спільної багатокористувацької роботи над справами консиліуму.
\end{enumerate}

Реалізація зазначених складових є передумовою побудови цілісної, масштабованої та керованої архітектури ЄСІКС, зниження витрат на розробку та супровід, забезпечення єдиних стандартів безпеки та користувацького досвіду.

\newpage
\section{Шаблони документів}
Централізована підсистема для розробки та зберігання шаблонів у форматах DOCX та XLSX з підтримкою динамічних плейсхолдерів.

\subsection{Вимоги до сутностей}
\subsubsection{Шаблон DOCX/XLSX}
Шаблон документу повинен бути у форматі OpenXML (ISO/IEC 29500) і підтримувати плейсхолдери таких типів: скалярні значення (текст), таблиці (завантаження з CSV), зображення (вставка логотипів та ілюстрацій).

\subsubsection{Набір змінних}
Реквізитна частина змінних включає векторні дані (таблиці), зображення та скалярні поля.

\subsection{Вимоги до процесів}
\subsubsection{Заповнення шаблону}
Процес заміни плейсхолдерів даними з вхідних параметрів (CLI-аргументи або API) з генерацією заповненого DOCX-файлу.

\subsubsection{Версіонування шаблонів}
Збереження та керування версіями шаблонів у централізованому сховищі.

Референсна імплементація надана в репозиторії \textbf{erpuno/docx}\footnote{\url{https://github.com/erpuno/docx}} (F\#-процесор шаблонів OpenXML).

\newpage
\section{Генерація документів}
Сервіс автоматичного формування документів у форматі PDF на основі затверджених шаблонів та даних із реєстрових систем ЄСІКС.

\subsection{Вимоги до сутностей}
\subsubsection{Заповнений документ}
Об’єкт, що містить заповнений DOCX/XLSX та мета-дані для конвертації.

\subsubsection{PDF-документ}
Фінальний артефакт у форматі PDF/A-3 (архівний) з вбудованими шрифтами та цифровим підписом.

\subsection{Вимоги до процесів}
\subsubsection{Генерація PDF}
Автоматична конвертація заповненого шаблону DOCX у PDF з збереженням форматування, таблиць та зображень.

\subsubsection{Накладання КЕП}
Автоматичне накладання кваліфікованого електронного підпису на згенерований PDF.

\subsubsection{Архівація}
Збереження PDF у довгострокове сховище з фіксацією хешу для цілісності.

\newpage
\section{Генерація бар- та штрих-кодів}
Підсистема генерації лінійних та матричних кодів (QR) для забезпечення простежуваності та верифікації паперових та електронних примірників документів.
Референсна імплементація надана в репозиторії \textbf{erpuno/barlix}\footnote{\url{https://github.com/erpuno/barlix}}
(BARLIX — генератор QR-кодів та лінійних штрих-кодів для Elixir з підтримкою Code39 та PNG-рендерингу).

\subsection{Вимоги до сутностей}
\subsubsection{Код}
Об’єкт, що містить тип (EAN-13, Code128, QR), дані (URL, ID документу, checksum) та зображення.

\subsubsection{Накладання}
Мета-дані позиції накладання на PDF/DOCX.

\subsection{Вимоги до процесів}
\subsubsection{Генерація коду}
Створення зображення коду за алгоритмами ISO/IEC 18004 (QR) та ISO/IEC 15417 (Code128).

\subsubsection{Накладання на документ}
Автоматичне вставлення коду в шаблон або PDF на заданій позиції.

\subsubsection{Верифікація}
Перевірка сканованого коду на відповідність оригіналу.

\newpage
\section{Розпізнавання Tesseract}
Інтегрований сервіс оптичного розпізнавання символів (OCR) для автоматизованого перетворення сканованих копій у текстовий формат.

\subsection{Вимоги до сутностей}
\subsubsection{Документ PDF}
Вхідний файл (сканований або текстовий PDF) з мета-даними (колір, статус).

\subsubsection{Результати OCR}
Текст з підтримкою української мови, векторні координати підсвічування, tsvector для повнотекстового пошуку.

\subsection{Вимоги до процесів}
\subsubsection{Інгестія документа}
Завантаження PDF та постановка в чергу на обробку (статус PENDING).

\subsubsection{Виконання OCR}
Розпізнавання за допомогою Tesseract OCR з українським словником (hunspell-uk, custom stop-words).

\subsubsection{Повнотекстовий пошук}
Індексація результатів у PostgreSQL з підтримкою префіксного пошуку та підсвічування.

\subsubsection{Отримання підсвічування}
Повернення координат термінів пошуку на сторінках документа.

Референсна імплементація надана в репозиторії \textbf{erpuno/ocr}\footnote{\url{https://github.com/erpuno/ocr}} (FastAPI + Tesseract + PostgreSQL).

\addtocontents{toc}{\protect\newpage}
\newpage
\section{Текстуально-голосові трансформації}
Компоненти для синтезу мовлення із тексту та розпізнавання голосових команд для інклюзивного доступу та автоматизації засідань.

\subsection{Вимоги до сутностей}
\subsubsection{Аудіо-потік}
Вхідний/вихідний аудіо-файл або потік у форматах WAV/OPUS.

\subsubsection{Текстовий фрагмент}
Сегмент тексту з тимінгами для S2T/T2S.

\subsection{Вимоги до процесів}
\subsubsection{Speech-to-Text (S2T)}
Розпізнавання мовлення (українська мова) у реальному часі або пакетному режимі.

\subsubsection{Text-to-Speech (T2S)}
Синтез мовлення з тексту з природним наголосом та інтонацією.

\subsubsection{Інтеграція з АРМ}
Автоматична транскрипція аудіо засідань та озвучення текстових рішень.

\newpage
\section{Система сканування}
Система сканування не повинна містити жодних ліцензій і бути повністю безкоштовною.
Протокол повинен підтримуватися TWAIN32. Leadtools, Scandit, Dynamsoft, IronOCR, PackageX, Tabscanner, VueScan, Mistral OCR та інші не підходять.
Система повинна бути написана на .NET і підтримувати Windows, Linux та Mac, та запускатися в Tray, як
фоновий процес, в середині якого повинен бути реалізований доступ до TWAIN пристроїв, а результати сканування
передаватися по WebSocket протоколу через цей фоновий процес в клієнтські застосунки які потребують джерел сканування.
Референсна імплементація надана в репозиторії \textbf{erpuno/scan}\footnote{\url{https://erpuno.github.io/scan}}.

\subsection{Вимоги до сутностей}
\subsubsection{Профіль сканера}
Налаштування TWAIN data-copy (Epson DS-530, Canon DR-C240, Kodak тощо) з параметрами duplex/autoscan.

\subsubsection{Сканований документ}
Багатосторінковий PDF з мета-даними (джерело, профіль, timestamp).

\subsection{Вимоги до процесів}
\subsubsection{Ініціація сканування}
Команда через WebSocket (наприклад, SCAN,DS-530,AUTOSCAN+AUTOFEED).

\subsubsection{Доставка результату}
Передача готового PDF у браузер/клієнт через WebSocket.

\newpage
\section{Багатокористувацький редактор}
Інструмент для спільної роботи суддів та секретарів над проєктами рішень у режимі реального часу на основі технології CRDT.

\subsection{Вимоги до сутностей}
\subsubsection{Редакційний документ}
Текстовий/структурований документ (справа консиліуму) з CRDT-структурами (текст, таблиці, коментарі).

\subsubsection{Операція редагування}
Атомарна зміна (insert/delete) з векторним годинником для конфліктного розв’язання.

\subsection{Вимоги до процесів}
\subsubsection{Синхронізація в реальному часі}
Одночасне редагування кількома користувачами без блокувань з мінімальним і контрольованим латенсі в VPN мережі.

\subsubsection{Історія змін}
Збереження повної історії операцій для undo/redo та аудиту.

\subsubsection{Фіксація версії}
Перехід у режим «тільки читання» після підписання КЕП.

\subsubsection{Транзакційність}
Транзакційність вимагається покомітно як в git, мінімальна одиниця редагування "параграф",
який повинен бути формально визначений в системі.

\newpage
\section*{Висновок}
Інфраструктура документообігу ЄСІКС створює умови для повного переходу до безпаперового судочинства, забезпечуючи автоматизацію рутинних операцій, безпеку зберігання та зручність колективної роботи над судовими справами.

\begin{thebibliography}{9}

\bibitem{ecoz:rehab} Міністерство охорони здоров’я України. \newblock Технічні вимоги Реабілітація 2.0. \newblock Версія 08.08.2024.
\bibitem{dstu:3973} ДСТУ 3973-2000. \newblock Система розроблення та постачання продукції на виробництво. Правила виконання науково-дослідних робіт. Загальні положення.
\bibitem{dstu:3974} ДСТУ 3974-2000. \newblock Система розроблення та постачання продукції на виробництво. Правила виконання науково-конструкторських робіт. Загальні положення.
\bibitem{dstu:42010} ДСТУ 42010:2018. \newblock Інженерія систем і програмних засобів. Опис архітектури.
\bibitem{law:info-protection} Закон України «Про захист інформації в інформаційно-телекомунікаційних системах».
\bibitem{law:personal-data} Закон України «Про захист персональних даних».
\bibitem{law:cybersecurity} Закон України «Про основні засади забезпечення кібербезпеки України».
\bibitem{owasp:top10} OWASP Foundation. \newblock OWASP Top 10: The Ten Most Critical Web Application Security Risks. \newblock \url{https://owasp.org/www-project-top-ten/}.
\bibitem{nist:sp800-57} NIST Special Publication 800-57 Part 1 Revision 5. \newblock Recommendation for Key Management.
\bibitem{iso:42010} ISO/IEC/IEEE 42010:2022. \newblock Systems and software engineering — Architecture description.
\bibitem{zachman} Zachman, J.A. \newblock A Framework for Information Systems Architecture. \newblock IBM Systems Journal, 1987.
\bibitem{esiks} Концепція ЄСІКС. \url{https://court.gov.ua/storage/portal/dsa/news/Kontseptsia%20ESIKS.pdf}. 2024
\bibitem{cbor} CBOR/COSE encoding/protocol. \url{https://tonpa.guru/stream/2024/2024-11-20%20CBOR%20COSE.htm}. 2025

\end{thebibliography}

\end{document}
